\documentclass[a4paper,11pt]{article}
\usepackage[T2A]{fontenc}
\usepackage[utf8]{inputenc}
\usepackage[russian]{babel}
\usepackage{amsmath,amssymb}
\usepackage{geometry}
\usepackage{microtype}
\usepackage{DejaVuSans}
\usepackage{needspace}
\renewcommand{\familydefault}{\sfdefault}
\usepackage{tabularx,booktabs,array}
\usepackage{enumitem}
\usepackage{parskip}
\usepackage{tikz}
\usepackage{xcolor}
\usepackage{hyperref}
\usepackage{fancyhdr}
\geometry{left=22mm,right=22mm,top=23mm,bottom=25mm}
\definecolor{accent}{HTML}{176E71}
\definecolor{navy}{HTML}{1F4055}
\definecolor{muted}{HTML}{617078}
\definecolor{error}{HTML}{A43836}
\hypersetup{colorlinks=true,linkcolor=accent,pdfauthor={Лёвин Артём Александрович},pdftitle={Ревью домашней работы — Екатерина Скелли}}
\setlength{\parindent}{0pt}
\setlength{\parskip}{4pt}
\setlist[enumerate]{leftmargin=*,itemsep=2pt,topsep=3pt}
\pagestyle{fancy}
\fancyhf{}
\renewcommand{\headrulewidth}{0pt}
\fancyfoot[C]{\footnotesize\color{muted} Лёвин Артём Александрович эксклюзивно для Екатерины Скелли}
\fancyfoot[R]{\footnotesize\color{muted}\thepage}
\newcommand{\field}[1]{\par\Needspace{4\baselineskip}\vspace{3pt}{\bfseries\color{navy}#1}\par\vspace{-2pt}}
\newcommand{\err}{\textcolor{error}{\bfseries Ошибка: }}
\newcommand{\idea}{\textcolor{accent}{\bfseries Ключевая идея: }}
\newcommand{\result}{\textcolor{accent}{\bfseries Итог: }}
\newcommand{\taskheading}[2]{\clearpage\phantomsection\label{task:#1}{\Large\bfseries\color{navy} Задание №\,#1. #2}\par\vspace{2pt}\textcolor{accent}{\rule{\linewidth}{0.6pt}}\vspace{7pt}}
\begin{document}
\thispagestyle{empty}
\vspace*{38mm}
\begin{center}
{\small\bfseries\color{accent} МАТЕМАТИКА}\par\vspace{12mm}
{\Huge\bfseries\color{navy} Ревью домашней работы}\par\vspace{7mm}
{\large Геометрия параллелограмма и системы счисления}\par\vspace{15mm}
{\large Екатерина Скелли}\par\vspace{5mm}
{\large 28 октября 2026 г.}\par\vspace{26mm}
{\normalsize Автор: Лёвин Артём Александрович}
\end{center}
\vfill
\begin{center}
\begin{tikzpicture}
\draw[accent!55,line width=1pt] (-4,0) .. controls (-1.6,0.65) and (1.5,-0.65) .. (4,0);
\end{tikzpicture}
\par\vspace{6mm}
{\footnotesize\color{muted} Лёвин Артём Александрович эксклюзивно для Екатерины Скелли}
\end{center}
\clearpage
\section*{Сводная таблица}
В таблице собраны задания, где нужно дополнить решение либо исправить доказательство. Верные вычисления внутри частично выполненных задач сохранены.
\par\vspace{5mm}
{\small
\setlength{\tabcolsep}{3pt}
\renewcommand{\arraystretch}{1.38}
\begin{tabularx}{\linewidth}{@{}p{12mm}p{26mm}X p{28mm}p{31mm}@{}}
\toprule
\textbf{№}&\textbf{Тема}&\textbf{Что исправить}&\textbf{Твой результат}&\textbf{Нужный результат}\\
\midrule
\hyperref[task:2]{2}&Углы&Найдены только два из четырёх углов.&$80^\circ$, $100^\circ$&$80^\circ$, $100^\circ$, $80^\circ$, $100^\circ$\\
\addlinespace[3pt]
\hyperref[task:5]{5}&Средние линии&Доказательство отсутствует.&Не записан&$ABCD$ — параллелограмм\\
\addlinespace[3pt]
\hyperref[task:6]{6}&Прямо\-угольник&Не обоснован прямой угол в построении.&Утверждение о прямоугольнике&Доказать прямой угол\\
\addlinespace[3pt]
\hyperref[task:7]{7}&Основание 16&Нет явной проверки допустимости цифр.&$59_{10}$&$59_{10}$; цифры допустимы\\
\addlinespace[3pt]
\hyperref[task:8]{8}&Основание 8&Пропущена обратная проверка.&$355_8$&$355_8$; проверка: $237_{10}$\\
\bottomrule
\end{tabularx}}
\vspace{9mm}
\textcolor{muted}{\small Каждая строка ведёт к подробному разбору соответствующего задания.}
\clearpage
\thispagestyle{empty}
\vspace*{\fill}
\begin{center}
{\Huge\bfseries\color{navy} Разбор ошибок}\par\vspace{5mm}
{\large\color{muted} Подробные пояснения и задания на закрепление}
\end{center}
\vspace*{\fill}

\taskheading{2}{Углы параллелограмма}
\field{Условие}
В параллелограмме $ABCD$ величины соседних углов $A$ и $B$ относятся как $4:5$. Найди все четыре угла.
\field{Твоё решение}
Использовано свойство суммы соседних углов:
\[
4x+5x=180^\circ,\qquad 9x=180^\circ,\qquad x=20^\circ.
\]
Найдены $\angle A=80^\circ$ и $\angle B=100^\circ$.
\field{Твой ответ}
$\angle A=80^\circ$, $\angle B=100^\circ$; остальные два угла не указаны.
\field{Где возникает неполнота}
\err Вычисление двух соседних углов выполнено правильно, но задание требует \emph{все четыре} угла. После нахождения $\angle B$ решение заканчивается раньше необходимого.
\field{Почему требуется ещё один шаг}
В параллелограмме противоположные углы равны:
\[
\angle A=\angle C,\qquad \angle B=\angle D.
\]
Без этих двух равенств итоговый ответ не содержит всех величин из условия.
\field{Ключевая идея}
\idea Сначала найди соседние углы, используя сумму $180^\circ$, затем восстанови два противоположных угла по их равенству.
\field{Исправление и полное решение}
\begin{enumerate}
\item Обозначим $\angle A=4x$, $\angle B=5x$.
\item Соседние углы параллелограмма составляют $180^\circ$: $4x+5x=180^\circ$.
\item Отсюда $x=20^\circ$; следовательно, $\angle A=80^\circ$, $\angle B=100^\circ$.
\item Противоположные углы равны: $\angle C=\angle A=80^\circ$, $\angle D=\angle B=100^\circ$.
\end{enumerate}
\field{Проверка}
$80^\circ+100^\circ=180^\circ$; противоположные углы попарно равны. Оба свойства выполнены.
\field{Итог}
\result $\angle A=\angle C=80^\circ$, $\angle B=\angle D=100^\circ$.
\field{Задача на закрепление}
В параллелограмме $KLMN$ углы $K$ и $L$ относятся как $3:7$. Найди все четыре угла.

\taskheading{5}{Параллелограмм из середин сторон}
\field{Условие}
В произвольном выпуклом четырёхугольнике $SKLE$ точки $A$, $B$, $C$, $D$ являются серединами сторон $SK$, $KL$, $LE$, $ES$ соответственно. Докажи, что $ABCD$ — параллелограмм.
\field{Твоё решение и ответ}
На листе отмечен номер задания и знак вопроса. Доказательство и вывод не записаны.
\field{Где возникает неполнота}
\err Для задачи на доказательство недостаточно назвать искомую фигуру. Нужно получить два признака параллельности противоположных сторон либо другой достаточный признак параллелограмма.
\field{Почему здесь помогает средняя линия}
\idea Отрезок, соединяющий середины двух сторон треугольника, параллелен третьей стороне и равен её половине. Проведём мысленно диагонали $SL$ и $KE$ исходного четырёхугольника. Они позволяют использовать эту теорему в четырёх треугольниках.
\field{Полное правильное доказательство}
\begin{enumerate}
\item Рассмотрим треугольник $SKL$. Точки $A$ и $B$ — середины $SK$ и $KL$. Поэтому $AB$ — средняя линия и $AB\parallel SL$.
\item Рассмотрим треугольник $SEL$. Точки $D$ и $C$ — середины $SE$ и $EL$. Поэтому $DC$ — средняя линия и $DC\parallel SL$.
\item Значит, $AB\parallel DC$: первая пара противоположных сторон параллельна.
\item В треугольнике $KLE$ точки $B$ и $C$ — середины $KL$ и $LE$, следовательно, $BC\parallel KE$.
\item В треугольнике $KSE$ точки $A$ и $D$ — середины $KS$ и $SE$, следовательно, $AD\parallel KE$.
\item Значит, $BC\parallel AD$: вторая пара противоположных сторон тоже параллельна.
\end{enumerate}
По определению параллелограмма четырёхугольник, у которого обе пары противоположных сторон параллельны, является параллелограммом. Следовательно, $ABCD$ — параллелограмм.
\field{Проверка}
Получены именно нужные пары: $AB\parallel CD$ и $BC\parallel AD$. Эти стороны принадлежат четырёхугольнику $ABCD$ и являются противоположными.
\field{Итог}
\result Доказано: $ABCD$ — параллелограмм. Достаточно применить теорему о средней линии четырежды.
\field{Задача на закрепление}
В выпуклом четырёхугольнике $PQRS$ точки $M$, $N$, $K$, $T$ — середины сторон $PQ$, $QR$, $RS$, $SP$ соответственно. Докажи, что $MNKT$ — параллелограмм.

\taskheading{6}{Как доказать, что параллелограмм — прямоугольник}
\field{Условие}
В параллелограмме $ABCD$ точка $M$ — середина стороны $AB$. Известно, что $MC=MD$. Докажи, что $ABCD$ — прямоугольник.
\field{Твоё решение}
Проведена вспомогательная конструкция: точка $E$ — середина $CD$, проведён отрезок $ME$. Верно замечено, что $\triangle MCD$ равнобедренный, поэтому медиана $ME$ к основанию $CD$ перпендикулярна $CD$. Затем сделана попытка обосновать $ME\parallel AD$ через равенство прямых углов.
\field{Твой ответ}
Сформулирован вывод о том, что $ABCD$ — прямоугольник, но обоснование содержит пропуск.
\field{Где именно возникает ошибка}
\err Последний надёжно доказанный факт — $ME\perp CD$. В следующем переходе угол между $CD$ и $DA$ фактически принимается равным $90^\circ$. Но именно перпендикулярность этих сторон требуется получить в задаче. Из рисунка её заключать нельзя.
\field{Почему так рассуждать нельзя}
На произвольном чертеже параллелограмма сторона $AD$ может быть наклонной по отношению к $CD$. Рисунок не доказывает, что между ними прямой угол. Если мы заранее назовём этот угол прямым, то используем требуемый вывод как уже известный факт. Получается логический круг.
\field{Ключевая идея}
\idea Твою конструкцию с точкой $E$ можно сохранить. Нужно \emph{независимо} доказать $ME\parallel AD$, пользуясь равенством и параллельностью половин противоположных сторон параллелограмма.
\field{Исправление от последнего правильного шага}
Так как $M$ и $E$ — середины сторон $AB$ и $CD$, имеем
\[
AM=\frac{AB}{2},\qquad ED=\frac{CD}{2}.
\]
В параллелограмме $AB=CD$, поэтому $AM=ED$. Кроме того, $AM\parallel ED$, поскольку $AB\parallel CD$. Следовательно, четырёхугольник $AMED$ — параллелограмм по признаку одной пары противоположных сторон, равных и параллельных. Значит, $ME\parallel AD$.
\newpage
{\small\color{muted} Задание №\,6. Продолжение доказательства}\par\vspace{3mm}
\field{Полное правильное доказательство}
\begin{enumerate}
\item По условию $MC=MD$. Значит, треугольник $MCD$ равнобедренный, его основание — $CD$.
\item Построим точку $E$ — середину $CD$ и проведём $ME$. В равнобедренном треугольнике медиана к основанию является высотой, поэтому $ME\perp CD$.
\item Поскольку $M$ — середина $AB$, $E$ — середина $CD$, а $AB=CD$ и $AB\parallel CD$, получаем $AM=ED$ и $AM\parallel ED$.
\item По признаку параллелограмма $AMED$ — параллелограмм. Значит, $ME\parallel AD$.
\item Из $ME\perp CD$ и $ME\parallel AD$ следует $AD\perp CD$. Итак, $\angle D=90^\circ$.
\item Параллелограмм с прямым углом — прямоугольник. Следовательно, $ABCD$ — прямоугольник.
\end{enumerate}
\field{Проверка}
Из $AD\perp CD$ и $AB\parallel CD$ также следует $AD\perp AB$. Значит, прямыми являются углы $D$ и $A$; противоположные им углы тоже прямые. Все четыре угла равны $90^\circ$.
\field{Итог}
\result Доказано: $ABCD$ — прямоугольник. Недостающим шагом было строгое доказательство $ME\parallel AD$.
\field{Задача на закрепление}
В параллелограмме $KLMN$ точка $P$ — середина стороны $KL$, причём $PM=PN$. Докажи, что $KLMN$ — прямоугольник.

\taskheading{7}{Допустимость шестнадцатеричной записи}
\field{Условие}
Переведи число $(3B)_{16}$ в десятичную систему и проверь, что запись допустима.
\field{Твоё решение}
Вычислено $3\cdot16^1+11\cdot16^0=59$.
\field{Твой ответ}
$(59)_{10}$ — числовой результат верен.
\field{Где возникает неполнота}
\err Выполнен перевод, но отдельно не записано обоснование того, что обе цифры допустимы для системы с основанием $16$.
\field{Почему это важно}
В позиционной системе с основанием $q$ разрешены цифры со значениями от $0$ до $q-1$. В шестнадцатеричной записи буквы $A,B,C,D,E,F$ обозначают числа от $10$ до $15$. Например, буква $G$ в обычной шестнадцатеричной записи недопустима, хотя её можно попытаться формально подставить в вычисление.
\field{Ключевая идея}
\idea Проверка допустимости — отдельное действие: установить значение каждой цифры и сравнить его с основанием.
\field{Исправление и полное решение}
\begin{enumerate}
\item Цифра $3$ обозначает число $3$, а $B$ — число $11$.
\item Для основания $16$ должно выполняться $0\le d<16$ для каждой цифры $d$.
\item Имеем $0\le3<16$ и $0\le11<16$: запись $(3B)_{16}$ допустима.
\item Переводим: $(3B)_{16}=3\cdot16+11=48+11=(59)_{10}$.
\end{enumerate}
\field{Проверка}
$59=3\cdot16+11$, где остаток $11$ соответствует цифре $B$. Обратное восстановление записи даёт $(3B)_{16}$.
\field{Итог}
\result Запись допустима; $(3B)_{16}=(59)_{10}$.
\field{Задача на закрепление}
Переведи число $(2E)_{16}$ в десятичную систему и отдельно проверь допустимость каждой его цифры.

\taskheading{8}{Обратная проверка перевода}
\field{Условие}
Переведи число $(237)_{10}$ в восьмеричную систему счисления. Выполни обратную проверку.
\field{Твоё решение}
Использовано последовательное деление на $8$, получено $(355)_8$.
\field{Твой ответ}
$(355)_8$ — перевод выполнен правильно.
\field{Где возникает неполнота}
\err После получения восьмеричной записи не выполнено действие, прямо указанное во второй части условия: обратная проверка переводом в десятичную систему.
\field{Почему проверка полезна}
При последовательном делении остатки записывают \emph{снизу вверх}. Если случайно поменять их порядок или ошибиться в частном, обратное вычисление по степеням основания обнаружит расхождение с исходным числом.
\field{Ключевая идея}
\idea Для обратной проверки умножь каждую цифру полученного числа на соответствующую степень восьмёрки и сложите результаты.
\field{Исправление и полное решение}
\begin{enumerate}
\item Делим число на $8$ с остатком:
\[
\begin{aligned}
237&=8\cdot29+5,\\
29&=8\cdot3+5,\\
3&=8\cdot0+3.
\end{aligned}
\]
\item Читаем остатки в обратном порядке: $3,5,5$. Получаем $(237)_{10}=(355)_8$.
\item Возвращаемся к основанию $10$:
\[
(355)_8=3\cdot8^2+5\cdot8^1+5\cdot8^0
=192+40+5=237.
\]
\end{enumerate}
\field{Проверка}
При обратном переводе получилось исходное число $237$. Значит, переход к восьмеричной системе подтверждён.
\field{Итог}
\result $(237)_{10}=(355)_8$, обратная проверка: $(355)_8=(237)_{10}$.
\field{Задача на закрепление}
Переведи число $(310)_{10}$ в восьмеричную систему счисления и выполни обратную проверку.
\end{document}